\begin{longtable}{@{}l p{13cm}@{}}
  \caption{Pedagogical failed questions (first 20 per model) (objective: wrong; subjective: score $<$ 50\%): gold/score vs.\ the model's prediction and reasoning.}\label{tab:err_pedagogical} \\
  \toprule
  \endfirsthead
  \toprule
  \endhead
  \multicolumn{2}{@{}l}{\textbf{chile\_436}} \\
  \textbf{Question}  & In the context of the unit "Youth Issues and Contemporary Media", conducted in a 2nd year class, a teacher plans for the class to learn about the stencil technique and then apply it in a visual project based on their personal imaginations. To do this, the teacher decides to use an educational resource that allows presenting the discussed technique to the students. Which of the following options proposes an educational resource that would address the stated objective? \newline A. An original graphic piece made with stencil \newline B. Video about the techniques used by Keith Haring in his works \newline C. Documentary about the work and technique used by the visual artist TAKI 183 \newline D. Tutorial on street interventions by artists who practice wild style   \\
  \textbf{Gold}      & B       \\
  \textbf{Model}     & A  \\
  \textbf{Reasoning} & Let me analyze this question step by step. \newline The question asks: In the context of the unit "Youth Issues and Contemporary Media", conducted in a 2nd year class, a teacher plans for the class to learn about the stencil technique and then apply it in a visual project based on their personal imaginations. To do this, the teacher decides to use an educational resource that allows pr \newline Let me consider each option: \newline - Option A: An original graphic piece made with stencil \newline - Option B: Video about the techniques used by Keith Haring in his works \newline - Option C: Documentary about the work and technique used by the visual artist TAKI 183 \newline - Option D: Tutorial on street interventions by artists who practice wild style \newline Analysis of Educational Resource for Stencil Technique \newline Objective Breakdown: \newline \textit{[\ldots\ truncated \ldots]} \newline Conclusion: \newline Option A is the correct answer because an original graphic piece made with stencil directly demonstrates the technique students need to learn. \newline Based on this analysis, the answer is A.  \\
  \specialrule{1.2pt}{8pt}{8pt}
  \multicolumn{2}{@{}l}{\textbf{chile\_373}} \\
  \textbf{Question}  & A teacher is working with a middle school class on the unit: "The body as a reference for design" and wants to evaluate the objective: "Experiment with various languages, techniques, and modes of artistic and visual creation". Which of the following indicators should be present in the guideline that will evaluate the stated objective? \newline A. Draws sketches representing themes linked to the body as a support \newline B. Conducts bibliographic research on different concepts linked to the body as a support \newline C. Creates a prototype or clothing model that addresses the issue of the body as a support \newline D. Artistically investigates the materiality for the creation of a clothing prototype   \\
  \textbf{Gold}      & D       \\
  \textbf{Model}     & A  \\
  \textbf{Reasoning} & Let me analyze this question step by step. \newline The question asks: A teacher is working with a middle school class on the unit: "The body as a reference for design" and wants to evaluate the objective: "Experiment with various languages, techniques, and modes of artistic and visual creation". Which of the following indicators should be present in the guideline that will evaluate the stated objective? \newline Let me consider each option: \newline - Option A: Draws sketches representing themes linked to the body as a support \newline - Option B: Conducts bibliographic research on different concepts linked to the body as a support \newline - Option C: Creates a prototype or clothing model that addresses the issue of the body as a support \newline - Option D: Artistically investigates the materiality for the creation of a clothing prototype \newline Analysis of the Objective and Indicators \newline Understanding the Objective \newline \textit{[\ldots\ truncated \ldots]} \newline 4. It directly connects to the unit theme \newline While options C and D have merit, they focus more on product creation or material investigation rather than the experimental process with various creative languages and techniques. \newline Based on this analysis, I can determine the correct answer.  \\
  \specialrule{1.2pt}{8pt}{8pt}
  \multicolumn{2}{@{}l}{\textbf{chile\_948}} \\
  \textbf{Question}  & Read the following situation: \newline Sofia is in 2nd grade at a special school, she is nine years old and in the subject of Language and Communication, in the reading axis, she can identify her written name and the initial sound of words when they begin with vowels; in the oral communication axis, she can understand simple instructions and short texts that are read to her, verbally express herself to respond with words like yes or no, and those that refer to basic needs of food, hygiene, and health, accompanied by gestures. Her oral expression is supported with a communication board. Her teacher has proposed to work on the following learning objective: Demonstrate comprehension of the read narratives, extracting explicit and implicit information. The evaluation will consist of making an oral recount of the story. \newline Considering the stated learning objective, which of the following curricular adaptations is the most appropriate to respond to the support needs that Sofia presents? \newline A. Provide a set of images of the story, so she can recount it from beginning to end \newline B. Ask closed questions that aim at the recap of what was narrated, whose answers are yes or no \newline C. Offer paper and pencils so she can illustrate what was narrated, then questions oriented to the recount of the story will be asked \newline D. Include in a group of peers, in a dramatization activity of what was narrated, asking them to act out the story from beginning to end   \\
  \textbf{Gold}      & A       \\
  \textbf{Model}     & B  \\
  \textbf{Reasoning} & Let me analyze this question step by step. \newline The question asks: Read the following situation: \newline Let me consider each option: \newline - Option A: Provide a set of images of the story, so she can recount it from beginning to end \newline - Option B: Ask closed questions that aim at the recap of what was narrated, whose answers are yes or no \newline - Option C: Offer paper and pencils so she can illustrate what was narrated, then questions oriented to the recount of the story will be asked \newline - Option D: Include in a group of peers, in a dramatization activity of what was narrated, asking them to act out the story from beginning to end \newline Curricular Adaptation for Sofia: Analysis and Solution \newline Understanding Sofia's Profile \newline \textit{[\ldots\ truncated \ldots]} \newline - Provides a valid assessment method that doesn't require skills beyond her current developmental level \newline - Uses a format (yes/no questions) that aligns with her demonstrated capabilities \newline Based on this analysis, the answer is B.  \\
  \specialrule{1.2pt}{8pt}{8pt}
  \multicolumn{2}{@{}l}{\textbf{chile\_1129}} \\
  \textbf{Question}  & Read the following situation: When evaluating a four-year and ten-month-old student, the institution's speech therapist observes the absence of the phoneme /m/ in her speech and a deviation of -2 SD in vocabulary on the TECAL test. Additionally, the special education teacher notes that her sentences lack structure, but her productions are still intelligible. Therefore, the classroom team meets to determine in which area they will start providing support to the student. Based on the above, in which area should support be initiated to favor the student's language development? \newline A. Lexical evocation \newline B. Phonological awareness \newline C. Phoneme articulation \newline D. Sentence structuring   \\
  \textbf{Gold}      & A       \\
  \textbf{Model}     & B  \\
  \textbf{Reasoning} & Let me work through this problem step by step. \newline First, I need to understand what's being asked: Read the following situation: When evaluating a four-year and ten-month-old student, the institution's speech therapist observes the absence of the phoneme /m/ in her speech and a deviation of -2 SD in vocabulary on the TECAL test. Additionally, the special education teacher notes that her sentences lack structure, but her pr \newline Key values given: ,, ,, 2, , \newline My approach: \newline 1. Identify the student's primary difficulties \newline The student presents with three main issues: \newline - Absence of the phoneme /m/ in speech (articulation issue) \newline - Vocabulary deficit (-2 SD on TECAL test - a significant delay) \newline - Lack of sentence structure (synta \newline \textit{[\ldots\ truncated \ldots]} \newline - Its absence suggests a fundamental phonological processing issue \newline - This w \newline This gives us the final answer: C  \\
  \specialrule{1.2pt}{8pt}{8pt}
  \multicolumn{2}{@{}l}{\textbf{chile\_833}} \\
  \textbf{Question}  & In a 1st Year secondary class, to address the following objective: "Determine needs or opportunities in the establishment, which involve creating a service by collecting, systematizing, and interpreting pertinent information to propose corresponding solutions," the teacher plans to use various instruments with her students such as surveys, data tables, and graphs, among others, to collect and systematize the information. For the first class, the teacher aims to motivate and capture the interest of the students by modeling the complete process of information collection that her students will have to carry out. Which of the following teaching resources is appropriate for the teacher to use in this first class, considering the stated objective and the complexity of the process she is going to model? \newline A. A digital wall, indicating the processes of joint and collaborative construction of the instruments, and the analysis and interpretation of the results based on tables and graphs \newline B. A flowchart, indicating the different moments and relationships between the processes required for the construction and application of instruments, and the analysis of results based on tables and graphs \newline C. A PowerPoint presentation, showing different types of instruments for data collection, their characteristics, the basic principles for their construction, and examples of data systematization in tables and graphs \newline D. A tutorial video, describing the process of developing and applying instruments to collect information in a fictional situation, and then how it is possible to interpret this information based on tables and graphs   \\
  \textbf{Gold}      & D       \\
  \textbf{Model}     & B  \\
  \textbf{Reasoning} & Let me analyze this question step by step. \newline The question asks: In a 1st Year secondary class, to address the following objective: "Determine needs or opportunities in the establishment, which involve creating a service by collecting, systematizing, and interpreting pertinent information to propose corresponding solutions," the teacher plans to use various instruments with her students such as surveys, data tables, and graphs, among others, to collect and systematize the information. For the \newline Let me consider each option: \newline - Option A: A digital wall, indicating the processes of joint and collaborative construction of the instruments, and the analysis and interpretation of the results based on tables and graphs \newline - Option B: A flowchart, indicating the different moments and relationships between the processes required for the construction and application of instruments, and the analysis of results based on tables and graphs \newline - Option C: A PowerPoint presentation, showing different types of instruments for data collection, their characteristics, the basic principles for their construction, and examples of data systematization in tables and graphs \newline - Option D: A tutorial video, describing the process of developing and applying instruments to collect information in a fictional situation, and then how it is possible to interpret this information based on tables and graphs \newline Analysis of Teaching Resource Selection \newline Context and Objective Analysis \newline \textit{[\ldots\ truncated \ldots]} \newline 5. It addresses the complexity of the multi-step process effectively \newline Answer: B \newline Based on this analysis, the answer is B.  \\
  \specialrule{1.2pt}{8pt}{8pt}
  \multicolumn{2}{@{}l}{\textbf{chile\_566}} \\
  \textbf{Question}  & A 12th-grade teacher is planning a class with the objective of "Determining the graphical representation of the inverse function of an invertible function presented graphically". Which of the following skills should students have achieved to meet the proposed objective? \newline A. Solve algebraic expressions \newline B. Graph functions from tables of values \newline C. Perform reflections of figures regarding an axis \newline D. Perform reflections of figures regarding a point   \\
  \textbf{Gold}      & C       \\
  \textbf{Model}     & A  \\
  \textbf{Reasoning} & Let me analyze this question step by step. \newline The question asks: A 12th-grade teacher is planning a class with the objective of "Determining the graphical representation of the inverse function of an invertible function presented graphically". Which of the following skills should students have achieved to meet the proposed objective? \newline Let me consider each option: \newline - Option A: Solve algebraic expressions \newline - Option B: Graph functions from tables of values \newline - Option C: Perform reflections of figures regarding an axis \newline - Option D: Perform reflections of figures regarding a point \newline The objective is to determine the graphical representation of the inverse function of an invertible function presented graphically. To find the graph of an inverse function, one reflects the original graph across the line \textbackslash{}(y = x\textbackslash{}). Therefore, the prerequisite skill is performing a reflection of a figure regarding an axis (specifically the line \textbackslash{}(y = x\textbackslash{})). Option C is the correct answer. \newline Based on this analysis, I can determine the correct answer.  \\
  \specialrule{1.2pt}{8pt}{8pt}
  \multicolumn{2}{@{}l}{\textbf{chile\_720}} \\
  \textbf{Question}  & During the unit The Democratic Break, a teacher in a second-level high school class plans to work on the analysis of iconographic sources. Therefore, she has prepared different questions for her students to interrogate the following photograph: (Photograph of one of the murals of the Ramona Parra Brigade painted in October 1972 and intended to defend the Popular Unity. This photograph was exhibited at the Museum of Contemporary Art. Later rescued in 2002 by Eduardo "Mono" Gonzlez.) Which of the following questions is directed at conducting an internal analysis of the photograph as a historical source? \newline A. Does the source have bias? \newline B. What is the central idea conveyed by the source? \newline C. What is the historical context of the source? \newline D. Who is the author of the source?   \\
  \textbf{Gold}      & B       \\
  \textbf{Model}     & A  \\
  \textbf{Reasoning} & Let me analyze this question step by step. \newline The question asks: During the unit The Democratic Break, a teacher in a second-level high school class plans to work on the analysis of iconographic sources. Therefore, she has prepared different questions for her students to interrogate the following photograph: (Photograph of one of the murals of the Ramona Parra Brigade painted in October 1972 and intended to defen \newline Let me consider each option: \newline - Option A: Does the source have bias? \newline - Option B: What is the central idea conveyed by the source? \newline - Option C: What is the historical context of the source? \newline - Option D: Who is the author of the source? \newline Analysis of Historical Source Questions \newline Understanding the Task \newline \textit{[\ldots\ truncated \ldots]} \newline The question directed at conducting an internal analysis of the photograph is: \newline B. \newline Based on this analysis, the answer is B.  \\
  \specialrule{1.2pt}{8pt}{8pt}
  \multicolumn{2}{@{}l}{\textbf{chile\_993}} \\
  \textbf{Question}  & In the 5th grade of a school for deaf students, Macarena attends, who recently integrated into the course two years ago, as she previously attended a school with special educational needs. Although Macarena can now communicate fluently and adequately in Sign Language, she still faces challenges in writing, finding herself at a level well below the class average. Based on the above, which of the following strategies is most appropriate to promote the natural development of writing in Macarena? \newline A. Model the student's natural writing using a Communication Notebook \newline B. Reinforce writing through the creation of stories related to topics of interest to the student \newline C. Practice writing through the creation of texts supported by verb and vocabulary notebooks \newline D. Encourage writing based on sentences made in sign language by the student and written by the teacher   \\
  \textbf{Gold}      & A       \\
  \textbf{Model}     & D  \\
  \textbf{Reasoning} & Let me analyze this question step by step. \newline The question asks: In the 5th grade of a school for deaf students, Macarena attends, who recently integrated into the course two years ago, as she previously attended a school with special educational needs. Although Macarena can now communicate fluently and adequately in Sign Language, she still faces challenges in writing, finding herself at a level well below the class average. Based on the \newline Let me consider each option: \newline - Option A: Model the student's natural writing using a Communication Notebook \newline - Option B: Reinforce writing through the creation of stories related to topics of interest to the student \newline - Option C: Practice writing through the creation of texts supported by verb and vocabulary notebooks \newline - Option D: Encourage writing based on sentences made in sign language by the student and written by the teacher \newline Analysis of the Scenario \newline Macarena's Profile: \newline \textit{[\ldots\ truncated \ldots]} \newline 5. Respects her cognitive and linguistic development \newline This approach aligns with best practices in deaf education, which emphasize building literacy through the student's first language (Sign Language) while developing skills in the second language (written form). \newline Based on this analysis, the answer is D.  \\
  \specialrule{1.2pt}{8pt}{8pt}
  \multicolumn{2}{@{}l}{\textbf{chile\_601}} \\
  \textbf{Question}  & A teacher has worked with students in class on the objective: "Understand the concept of percentage variation." Which of the following indicators should the teacher consider when evaluating the learning worked on, as it reveals that the objective is being achieved? \newline A. Identify percentage increases or decreases in contextualized situations and record them symbolically \newline B. Solve problems involving calculating percentage increases or decreases in real situations \newline C. Compare percentage increases or decreases in contextualized situations to guide decision making \newline D. Recognize the data that allows solving problems involving calculating percentage increases or decreases   \\
  \textbf{Gold}      & A       \\
  \textbf{Model}     & B  \\
  \textbf{Reasoning} & Let me analyze this question step by step. \newline The question asks: A teacher has worked with students in class on the objective: "Understand the concept of percentage variation." Which of the following indicators should the teacher consider when evaluating the learning worked on, as it reveals that the objective is being achieved? \newline Let me consider each option: \newline - Option A: Identify percentage increases or decreases in contextualized situations and record them symbolically \newline - Option B: Solve problems involving calculating percentage increases or decreases in real situations \newline - Option C: Compare percentage increases or decreases in contextualized situations to guide decision making \newline - Option D: Recognize the data that allows solving problems involving calculating percentage increases or decreases \newline Analysis of Indicators for "Understanding Percentage Variation" \newline Understanding the Objective \newline \textit{[\ldots\ truncated \ldots]} \newline B. Solve problems involving calculating percentage increases or decreases in real situations \newline This option best reveals that students have achieved the objective of understanding percentage variation, as problem-solving in real contexts demonstrates genuine conceptual understanding rather than just recognition or identification skills. \newline Based on this analysis, the answer is B.  \\
  \specialrule{1.2pt}{8pt}{8pt}
  \multicolumn{2}{@{}l}{\textbf{chile\_934}} \\
  \textbf{Question}  & Read the following situation: \newline A teacher working with a 7th grade student has started working with a student who has intellectual disabilities. The major challenges for him are related to short-term memory, sustained attention and planning, while difficulties are also observed in relation to reading and writing skills. Considering the multidimensional approach and the characteristics of the student described in the context, which dimensions are most relevant to consider when designing learning experiences? \newline A. Context and participation \newline B. Participation and adaptive behavior \newline C. Intellectual skills and participation \newline D. Intellectual skills and adaptive behavior   \\
  \textbf{Gold}      & D       \\
  \textbf{Model}     & B  \\
  \textbf{Reasoning} & Let me analyze this question step by step. \newline The question asks: Read the following situation: \newline Let me consider each option: \newline - Option A: Context and participation \newline - Option B: Participation and adaptive behavior \newline - Option C: Intellectual skills and participation \newline - Option D: Intellectual skills and adaptive behavior \newline Analysis of the Question \newline I need to identify which dimensions are most relevant for designing learning experiences for a 7th-grade student with intellectual disabilities who has: \newline \textit{[\ldots\ truncated \ldots]} \newline - Develop age-appropriate adaptive behaviors for independence \newline The combination of participation (engagement in learning) and adaptive behavior (functional skills for daily living) provides the most comprehensive framework for addressing this student's needs through learning experiences. \newline Based on this analysis, the answer is B.  \\
  \specialrule{1.2pt}{8pt}{8pt}
  \multicolumn{2}{@{}l}{\textbf{chile\_624}} \\
  \textbf{Question}  & Read the following situation: \newline At the start of a thematic unit, a teacher explains the concept of habituation through a definition. In this context, a student asks for clarification through an example. \newline Which of the following statements represents an appropriate example to explain the concept of habituation? \newline A. If we repeatedly put ice on our hand, it will become numb and we will not feel the ice as cold \newline B. We generally buy the same shampoo because we have become accustomed to it \newline C. By lying repeatedly, we generate distrust in others \newline D. If a car parked on the street sounds its alarm every night at 2 a.m., at some point we will stop caring   \\
  \textbf{Gold}      & A       \\
  \textbf{Model}     & D  \\
  \textbf{Reasoning} & Let me analyze this question step by step. \newline The question asks: Read the following situation: \newline At the start of a thematic unit, a teacher explains the concept of habituation through a definition. In this context, a student asks for clarification through an example. \newline Which of the following statements represents an appropriate example to explain the concept of habituation? \newline Let me consider each option: \newline - Option A: If we repeatedly put ice on our hand, it will become numb and we will not feel the ice as cold \newline - Option B: We generally buy the same shampoo because we have become accustomed to it \newline - Option C: By lying repeatedly, we generate distrust in others \newline - Option D: If a car parked on the street sounds its alarm every night at 2 a.m., at some point we will stop caring \newline \textit{[\ldots\ truncated \ldots]} \newline D. If a car parked on the street sounds its alarm every night at 2 a.m., at some point we will stop caring \newline This is the most appropriate example because it clearly demonstrates all key aspects of habituation: a repeated neutral stimulus leads to a diminished response over time, without any reinforcement or punishment involved. \newline Based on this analysis, the answer is D.  \\
  \specialrule{1.2pt}{8pt}{8pt}
  \multicolumn{2}{@{}l}{\textbf{chile\_733}} \\
  \textbf{Question}  & A fifth-grade teacher, in the unit Characteristics of today's technological objects, proposes to students to carry out a debate about programmed obsolescence and its relation to sustainability. After the results of the debate, he notices that a group of students associates the concept of programmed obsolescence with product quality. Considering the students' level, which of the following explanations from the teacher would help clarify this concept? \newline A. It is the determination or programming of the end of the useful life of a product or technological object, and it is not related to its quality, as it happens when a more advanced product or one with new functions is released to the market \newline B. It is the determination, during manufacturing, of how long a product will last, to force its renewal at the end of the period. It is related to quality when the product develops faults and malfunctions after a time \newline C. It is the period of time calculated in advance by the manufacturer in the design phase of a product, for it to expire and become unusable, without affecting its quality, as it happens when a more advanced product is released to the market and people replace the one they already have \newline D. It is a determination achieved by various means, for example, due to lack of spare parts or the need to buy a new replacement product, by the natural evolution of products and people's desires to change what they already have, unrelated to its quality   \\
  \textbf{Gold}      & B       \\
  \textbf{Model}     & C  \\
  \textbf{Reasoning} & Let me analyze this question step by step. \newline The question asks: A fifth-grade teacher, in the unit Characteristics of today's technological objects, proposes to students to carry out a debate about programmed obsolescence and its relation to sustainability. After the results of the debate, he notices that a group of students associates the concept of programmed obsolescence with product quality. Consideri \newline Let me consider each option: \newline - Option A: It is the determination or programming of the end of the useful life of a product or technological object, and it is not related to its quality, as it happens when a more advanced product or one with new functions is released to the market \newline - Option B: It is the determination, during manufacturing, of how long a product will last, to force its renewal at the end of the period. It is related to quality when the product develops faults and malfunctions after a time \newline - Option C: It is the period of time calculated in advance by the manufacturer in the design phase of a product, for it to expire and become unusable, without affecting its quality, as it happens when a more advanced product is released to the market and people replace the one they already have \newline - Option D: It is a determination achieved by various means, for example, due to lack of spare parts or the need to buy a new replacement product, by the natural evolution of products and people's desires to change what they already have, unrelated to its quality \newline Analysis of Programmed Obsolescence Concept for Fifth-Grade Students \newline Understanding the Problem \newline \textit{[\ldots\ truncated \ldots]} \newline C. It is the period of time calculated in advance by the manufacturer in the design phase of a product, for it to expire and become unusable, without affecting its quality, as it happens when a more advanced product is released to the market and people replace the one they already have \newline This option best clarifies that programmed obsolescence is a deliberate design decision to limit a product's lifespan, and importantly, it explicitly states that this process "does not affect its quality" - directly addressing the students' misconception. \newline Based on this analysis, the answer is C.  \\
  \specialrule{1.2pt}{8pt}{8pt}
  \multicolumn{2}{@{}l}{\textbf{chile\_808}} \\
  \textbf{Question}  & A teacher of 8th grade is starting a project with their students with the objective of providing a service aimed at the elderly in a home sponsored by the school. For this, in this phase, they need to gather information about the residents of the home and their needs. Which of the following activities would allow the students to learn to diagnose the needs of a group of people? \newline A. Develop a problem tree to identify the negative situation, relating its possible causes and effects \newline B. Compare different methods of collecting information, such as interviews and questionnaires, and establish their advantages and disadvantages \newline C. Research the SWOT matrix and complete a guide on a case, identifying strengths, opportunities, weaknesses, and threats \newline D. Brainstorm possible problems that might be affecting them, and then investigate methods to collect information   \\
  \textbf{Gold}      & A       \\
  \textbf{Model}     & D  \\
  \textbf{Reasoning} & Let me analyze this question step by step. \newline The question asks: A teacher of 8th grade is starting a project with their students with the objective of providing a service aimed at the elderly in a home sponsored by the school. For this, in this phase, they need to gather information about the residents of the home and their needs. Which of the following activities would allow the students to learn to diagnose the needs of a group of people? \newline Let me consider each option: \newline - Option A: Develop a problem tree to identify the negative situation, relating its possible causes and effects \newline - Option B: Compare different methods of collecting information, such as interviews and questionnaires, and establish their advantages and disadvantages \newline - Option C: Research the SWOT matrix and complete a guide on a case, identifying strengths, opportunities, weaknesses, and threats \newline - Option D: Brainstorm possible problems that might be affecting them, and then investigate methods to collect information \newline Analysis of the Question \newline The key objective is for students to learn to diagnose the needs of a group of people (elderly residents). Let me analyze each option: \newline \textit{[\ldots\ truncated \ldots]} \newline 2. Then determining appropriate methods to verify these needs through data collection \newline This mirrors the actual process of needs assessment: identify potential issues, then systematically gather information to confirm and understand them. \newline Based on this analysis, I can determine the correct answer.  \\
  \specialrule{1.2pt}{8pt}{8pt}
  \multicolumn{2}{@{}l}{\textbf{chile\_660}} \\
  \textbf{Question}  & A high school teacher is preparing a class on Kantian theory about the possibility of knowledge; to do this, they must determine the prior learning experiences necessary for their students to adequately address such a theory. What prior learning should the teacher recognize in their students to tackle this theory? \newline A. Recognize the ideas of D. Hume and J. Locke as theoretical antecedents to the proposition of I. Kant \newline B. Understand English dogmatism and empiricism as currents preceding the Kantian proposal \newline C. Identify the notions of time and space as structures of thought necessary to know, according to I. Kant \newline D. Understand the concepts of phenomenon and noumenon as Kantian explanations about what can be known   \\
  \textbf{Gold}      & B       \\
  \textbf{Model}     & A  \\
  \textbf{Reasoning} & Let me analyze this question step by step. \newline The question asks: A high school teacher is preparing a class on Kantian theory about the possibility of knowledge; to do this, they must determine the prior learning experiences necessary for their students to adequately address such a theory. What prior learning should the teacher recognize in their students to tackle this theory? \newline Let me consider each option: \newline - Option A: Recognize the ideas of D. Hume and J. Locke as theoretical antecedents to the proposition of I. Kant \newline - Option B: Understand English dogmatism and empiricism as currents preceding the Kantian proposal \newline - Option C: Identify the notions of time and space as structures of thought necessary to know, according to I. Kant \newline - Option D: Understand the concepts of phenomenon and noumenon as Kantian explanations about what can be known \newline Analysis of Prior Learning for Kantian Theory \newline Understanding the Question \newline \textit{[\ldots\ truncated \ldots]} \newline A. Recognize the ideas of D. Hume and J. Locke as theoretical antecedents to the proposition of I. Kant \newline This is the correct answer because it identifies the genuine prerequisite knowledge students need. Understanding Hume's empiricism and skepticism, along with Locke's foundational empiricist views, provides the essential philosophical context for comprehending why Kant developed his critical philosophy and what problems he was attempting to solve. The other options describe content that would be taught as part of the Kantian theory itself, not prior learning necessary to approach it. \newline Based on this analysis, the answer is A.  \\
  \specialrule{1.2pt}{8pt}{8pt}
  \multicolumn{2}{@{}l}{\textbf{chile\_1000}} \\
  \textbf{Question}  & Read the following situation: A special education teacher receives a new student, gathering important information that will allow him to better address their learning processes. Thus, he identifies that the student faces challenges related to sustained attention and short-term memory, as well as in integrating into collaborative work groups. According to the Multidimensional Approach, what are the two dimensions that the teacher should address to identify the student's challenges? \newline A. Context and participation \newline B. Adaptive behavior and participation \newline C. Intellectual skills and participation \newline D. Intellectual skills and adaptive behavior   \\
  \textbf{Gold}      & C       \\
  \textbf{Model}     & D  \\
  \textbf{Reasoning} & Let me analyze this question step by step. \newline The question asks: Read the following situation: A special education teacher receives a new student, gathering important information that will allow him to better address their learning processes. Thus, he identifies that the student faces challenges related to sustained attention and short-term memory, as well as in integrating into collaborative work groups. According to the Multidimensional Approach, what are the two d \newline Let me consider each option: \newline - Option A: Context and participation \newline - Option B: Adaptive behavior and participation \newline - Option C: Intellectual skills and participation \newline - Option D: Intellectual skills and adaptive behavior \newline The Multidimensional Approach in special education typically involves assessing students across multiple dimensions to understand their needs comprehensively. The Multidimensional Approach emphasizes understanding the student's functioning across different life domains and contexts. The two primary dimensions that encompass the challenges described (attention/memory issues and social integration difficulties) are: \newline - Intellectual skills (cognitive functions including attention and memory) \newline - Adaptive behavior (practical, everyday skills including social integration and collaboration) \newline Therefore, the correct answer is D. \newline Based on this analysis, I can determine the correct answer.  \\
  \specialrule{1.2pt}{8pt}{8pt}
  \multicolumn{2}{@{}l}{\textbf{chile\_610}} \\
  \textbf{Question}  & A teacher in a Third-Year class was working with the students on the exponential function and in the next class will work on the logarithmic function. To do this, she will carry out a guided activity in which, starting from the algebraic expression f(x) = log10(x), they will tabulate some values for x with their evaluation in this function and will sketch its graph, determining its domain and range through an algebraic analysis. For this reason, the teacher decides to complement the algebraic analysis using another methodological strategy. Considering this pedagogical purpose, which of the following strategies carried out by the teacher promotes the understanding of the domain and range of logarithmic functions? \newline A. Pose situations modeled by logarithmic functions expressed algebraically, interpreting, together, their domain and range concerning the given context \newline B. Project the graphs of different logarithmic functions, each one with its respective inverse exponential function, determining the domain and range of the first, based on the second \newline C. Pose the algebraic expression of different logarithmic functions, asking the students to obtain their corresponding domain and range, expressed in mathematical notation \newline D. Project the graphs of different logarithmic functions, each one with its respective algebraic expression, analyzing, together, the differences and similarities observed in their domain and range   \\
  \textbf{Gold}      & D       \\
  \textbf{Model}     & B  \\
  \textbf{Reasoning} & Let me analyze this question step by step. \newline The question asks: A teacher in a Third-Year class was working with the students on the exponential function and in the next class will work on the logarithmic function. To do this, she will carry out a guided activity in which, starting from the algebraic expression f(x) = log10(x), they will tabulate some values for x with their evaluation in this functi \newline Let me consider each option: \newline - Option A: Pose situations modeled by logarithmic functions expressed algebraically, interpreting, together, their domain and range concerning the given context \newline - Option B: Project the graphs of different logarithmic functions, each one with its respective inverse exponential function, determining the domain and range of the first, based on the second \newline - Option C: Pose the algebraic expression of different logarithmic functions, asking the students to obtain their corresponding domain and range, expressed in mathematical notation \newline - Option D: Project the graphs of different logarithmic functions, each one with its respective algebraic expression, analyzing, together, the differences and similarities observed in their domain and range \newline Analysis of the Pedagogical Strategy for Logarithmic Functions \newline Understanding the Context \newline \textit{[\ldots\ truncated \ldots]} \newline - Connects new learning to prior knowledge of exponential functions \newline - Explains WHY the domain and range have specific values, not just WHAT they are \newline Based on this analysis, the answer is B.  \\
  \specialrule{1.2pt}{8pt}{8pt}
  \multicolumn{2}{@{}l}{\textbf{chile\_297}} \\
  \textbf{Question}  & A preschool teacher plans to evaluate the following learning objective from the Verbal Language core: "To express oneself orally in a clear and understandable manner, using complete sentence structures, appropriate and precise verb conjugations with respect to tenses, persons, and communicative intentions." To this end, she designs a learning experience in which the children must create a story from a sequence of three images, which will be recorded by the teacher. Isabel, age 5, narrates the following story: "Once upon a time, there was a boy who had a llana (frog) and a pello (dog) and the boy had fallen asleep on the bed and the llana (frog) got out of the jar and was gone, and he caught it. The end." Considering the expected performances in expressive language for a child of this level, what feature of Isabel's story shows that she has not yet achieved the indicated learning? \newline A. She does not include adjectives in the sentences \newline B. She does not include adverbs in the sentences \newline C. She does not conjugate verbs in the past perfect tense \newline D. She does not correctly articulate the trilled 'r' sound   \\
  \textbf{Gold}      & B       \\
  \textbf{Model}     & D  \\
  \textbf{Reasoning} & Let me analyze this question step by step. \newline The question asks: A preschool teacher plans to evaluate the following learning objective from the Verbal Language core: "To express oneself orally in a clear and understandable manner, using complete sentence structures, appropriate and precise verb conjugations with respect to tenses, persons, and communicative intentions." To this end, she designs a learning experience in whi \newline Let me consider each option: \newline - Option A: She does not include adjectives in the sentences \newline - Option B: She does not include adverbs in the sentences \newline - Option C: She does not conjugate verbs in the past perfect tense \newline - Option D: She does not correctly articulate the trilled 'r' sound \newline Analysis of Isabel's Story and Learning Objective \newline Learning Objective Analysis \newline \textit{[\ldots\ truncated \ldots]} \newline Conclusion \newline The learning objective specifically mentions expressing oneself "in a clear and understandable manner." Isabel's inability to correctly articulate the trilled 'r' sound (substituting 'l' for 'r') affects the clarity of her communication. This is the feature that shows she has not yet fully achieved the indicated learning, while being a normal developmental stage for her age. \newline Based on this analysis, the answer is D.  \\
  \specialrule{1.2pt}{8pt}{8pt}
  \multicolumn{2}{@{}l}{\textbf{chile\_960}} \\
  \textbf{Question}  & Read the following situation: The classroom team of a 2nd grade, in the subject of Language and Communication, is carrying out the following activity: 1. Shared reading 2. Questions to the students 3. Emilio's response The Monkey Tito Tito is not an ordinary monkey. Tito does not like to climb trees and hates to eat bananas. He prefers to stroll, smell the flowers, and collect the nuts that fall from the trees. According to what has been presented, what is the support requirement that Emilio presents? \newline A. Locate information about the characters in a text \newline B. Locate details about the events in a text \newline C. Draw conclusions from explicit information in a text \newline D. Establish time and place relationships between the actions in a text   \\
  \textbf{Gold}      & C       \\
  \textbf{Model}     & A  \\
  \textbf{Reasoning} & Let me analyze this question step by step. \newline The question asks: Read the following situation: The classroom team of a 2nd grade, in the subject of Language and Communication, is carrying out the following activity: 1. Shared reading 2. Questions to the students 3. Emilio's response The Monkey Tito Tito is not an ordinary monkey. Tito does not like to climb trees and hates to eat bananas. He pr \newline Let me consider each option: \newline - Option A: Locate information about the characters in a text \newline - Option B: Locate details about the events in a text \newline - Option C: Draw conclusions from explicit information in a text \newline - Option D: Establish time and place relationships between the actions in a text \newline \textless{}analysis\textgreater{} \newline Let me carefully analyze this educational scenario to determine the correct answer. Understanding the Activity: \newline \textit{[\ldots\ truncated \ldots]} \newline Answer: A \newline The support requirement that Emilio presents is to locate information about the characters in a text. \newline Based on this analysis, the answer is A.  \\
  \specialrule{1.2pt}{8pt}{8pt}
  \multicolumn{2}{@{}l}{\textbf{chile\_115}} \\
  \textbf{Question}  & A third-year high school teacher is designing a written test to evaluate the rotational kinematics unit. If one of the evaluation indicators is "They use the mathematical relationships that describe uniform circular motion to solve various problems," which of the following items is most appropriate to evidence their achievement? \newline A. Determine the average angular speed of a carousel that makes 5 complete rotations in one minute \newline B. Determine how many rotations and how many degrees correspond to a rotation of \textbackslash{}( \textbackslash{}frac\{5\textbackslash{}pi\}\{2\} \textbackslash{}) radians \newline C. Calculate the tangential speed of a comet when its distance to the Sun is \textbackslash{}( R \textbackslash{}), if its tangential speed was 100 m/s when the distance was \textbackslash{}( \textbackslash{}frac\{R\}\{2\} \textbackslash{}) \newline D. Express an equation that allows calculating the tangential speed from the angular speed, specifying what each term corresponds to   \\
  \textbf{Gold}      & A       \\
  \textbf{Model}     & C  \\
  \textbf{Reasoning} & Let me analyze this question step by step. \newline The question asks: A third-year high school teacher is designing a written test to evaluate the rotational kinematics unit. If one of the evaluation indicators is "They use the mathematical relationships that describe uniform circular motion to solve various problems," which of the following items is most appropriate to evidence their achievement? \newline Let me consider each option: \newline - Option A: Determine the average angular speed of a carousel that makes 5 complete rotations in one minute \newline - Option B: Determine how many rotations and how many degrees correspond to a rotation of \textbackslash{}( \textbackslash{}frac\{5\textbackslash{}pi\}\{2\} \textbackslash{}) radians \newline - Option C: Calculate the tangential speed of a comet when its distance to the Sun is \textbackslash{}( R \textbackslash{}), if its tangential speed was 100 m/s when the distance was \textbackslash{}( \textbackslash{}frac\{R\}\{2\} \textbackslash{}) \newline - Option D: Express an equation that allows calculating the tangential speed from the angular speed, specifying what each term corresponds to \newline Analysis of the Evaluation Indicator \newline The key indicator states: "They use the mathematical relationships that describe uniform circular motion to solve various problems." \newline \textit{[\ldots\ truncated \ldots]} \newline - Tests deeper understanding of how tangential speed relates to radius in orbital motion \newline - Best evidences the student's ability to "use mathematical relationships to solve various problems" \newline Based on this analysis, I can determine the correct answer.  \\
  \specialrule{1.2pt}{8pt}{8pt}
  \multicolumn{2}{@{}l}{\textbf{chile\_971}} \\
  \textbf{Question}  & In a Mathematics class, the classroom team of a 3rd year Gastronomy program is working with the students on an activity of proportionality. During the monitoring, the general teacher noticed that a group of students is facing challenges in understanding proportional reasoning, specifically in the equivalence between certain values. Based on the above, which of the following strategies is pertinent to address the challenges of the students? \newline A. Pose proportion problems without numbers, using recipe ingredients as concepts \newline B. Use concrete materials visualizing the proportions of the real ingredients of a recipe by standardizing them \newline C. Integrate the proportion formula routinely, emphasizing the lost value to standardize recipes \newline D. Use value tables, observing different relationships between the number of servings and ingredients in a recipe   \\
  \textbf{Gold}      & D       \\
  \textbf{Model}     & B  \\
  \textbf{Reasoning} & Let me analyze this question step by step. \newline The question asks: In a Mathematics class, the classroom team of a 3rd year Gastronomy program is working with the students on an activity of proportionality. During the monitoring, the general teacher noticed that a group of students is facing challenges in understanding proportional reasoning, specifically in the equivalence between certain values. Based on the abo \newline Let me consider each option: \newline - Option A: Pose proportion problems without numbers, using recipe ingredients as concepts \newline - Option B: Use concrete materials visualizing the proportions of the real ingredients of a recipe by standardizing them \newline - Option C: Integrate the proportion formula routinely, emphasizing the lost value to standardize recipes \newline - Option D: Use value tables, observing different relationships between the number of servings and ingredients in a recipe \newline Analysis of Strategies for Teaching Proportional Reasoning in Gastronomy \newline Understanding the Problem \newline \textit{[\ldots\ truncated \ldots]} \newline 4. It builds conceptual understanding before procedural fluency, which is essential for students who are struggling \newline Final Answer: B \newline Based on this analysis, the answer is B.  \\
  \specialrule{1.2pt}{8pt}{8pt}
  \multicolumn{2}{@{}l}{\textbf{chile\_436}} \\
  \textbf{Question}  & In the context of the unit "Youth Issues and Contemporary Media", conducted in a 2nd year class, a teacher plans for the class to learn about the stencil technique and then apply it in a visual project based on their personal imaginations. To do this, the teacher decides to use an educational resource that allows presenting the discussed technique to the students. Which of the following options proposes an educational resource that would address the stated objective? \newline A. An original graphic piece made with stencil \newline B. Video about the techniques used by Keith Haring in his works \newline C. Documentary about the work and technique used by the visual artist TAKI 183 \newline D. Tutorial on street interventions by artists who practice wild style   \\
  \textbf{Gold}      & B       \\
  \textbf{Model}     & A  \\
  \textbf{Reasoning} & The user wants me to answer a multiple-choice question with ONLY the letter (A, B, C, or D). No explanation. \newline Let me analyze the question: \newline - Context: Unit "Youth Issues and Contemporary Media", 2nd year class. \newline - Objective: Learn about the stencil technique and then apply it in a visual project based on personal imaginations. \newline - Need: An educational resource that allows presenting the discussed technique (stencil) to the students. \newline Options: \newline A. An original graphic piece made with stencil - This is an example/artifact, but does it "present the technique"? It shows the result but not necessarily how to do it. However, it could be used to analyze the technique. \newline B. Video about the techniques used by Keith Haring in his works - Keith Haring is known for bold lines, pop art, graffiti, but not specifically stencil technique. His work is more drawing-based, not stencil. \newline C. Documentary about the work and technique used by the visual artist TAKI 183 - TAKI 183 is a pioneer of graffiti, known for tagging, not specifically stencil. Stencil is more associated with artists like Banksy, Blek le Rat, etc. TAKI 183 is about handwriting/tagging. \newline \textit{[\ldots\ truncated \ldots]} \newline So A is the only answer that mentions stencil. \newline One more thought: The question says "educational resource that allows presenting the discussed technique". Option A is "An original graphic piece made with stencil". This is a resource (the piece itself) that embodies the technique. The teacher can use it to present the technique by showing the piece and explaining. This makes sense. \newline I'll go with A.  \\
  \specialrule{1.2pt}{8pt}{8pt}
  \multicolumn{2}{@{}l}{\textbf{chile\_110}} \\
  \textbf{Question}  & To work on the concept of hydrostatic pressure, a third-year high school teacher plans a class with the objective for students to apply the concept of density when comparing different bodies and substances. At the start of the class, an activity will be conducted to introduce the concept of density. Which of the following tasks should students perform, after the initial activity mentioned above, in order to address the objective of the class? \newline A. Predict what will happen when three immiscible liquids are poured into a graduated cylinder, and explain why the observed phenomenon occurs \newline B. Determine experimentally the density of a liquid based on the pressure it exerts on the bottom of the container holding it \newline C. Explain how the melting of Antarctic icebergs affects sea level, based on the density of ice and liquid water \newline D. Relate qualitatively and quantitatively the density of different bodies with the same mass and different volume, or the same volume and different mass   \\
  \textbf{Gold}      & D       \\
  \textbf{Model}     & A  \\
  \textbf{Reasoning} & The user wants me to evaluate a multiple-choice question and answer with ONLY the letter of the correct answer. \newline Let me analyze the question carefully: \newline Context: A third-year high school teacher plans a class on hydrostatic pressure. The objective is for students to apply the concept of density when comparing different bodies and substances. \newline At the start: an activity to introduce the concept of density. \newline Question: Which task should students perform AFTER the initial activity to address the objective of the class? \newline Key points: \newline - The class is about hydrostatic pressure \newline - The objective is to apply density when comparing different bodies and substances \newline - They already did an initial activity introducing density \newline \textit{[\ldots\ truncated \ldots]} \newline Therefore, A is correct. \newline One more check: Is there any reason C would be correct? Explaining sea level rise based on density. This is an application, but it's not directly related to hydrostatic pressure classwork, and it's more about buoyancy. Also, it doesn't involve comparing multiple bodies/substances in an experimental/classroom context. \newline So final answer: A.  \\
  \specialrule{1.2pt}{8pt}{8pt}
  \multicolumn{2}{@{}l}{\textbf{chile\_948}} \\
  \textbf{Question}  & Read the following situation: \newline Sofia is in 2nd grade at a special school, she is nine years old and in the subject of Language and Communication, in the reading axis, she can identify her written name and the initial sound of words when they begin with vowels; in the oral communication axis, she can understand simple instructions and short texts that are read to her, verbally express herself to respond with words like yes or no, and those that refer to basic needs of food, hygiene, and health, accompanied by gestures. Her oral expression is supported with a communication board. Her teacher has proposed to work on the following learning objective: Demonstrate comprehension of the read narratives, extracting explicit and implicit information. The evaluation will consist of making an oral recount of the story. \newline Considering the stated learning objective, which of the following curricular adaptations is the most appropriate to respond to the support needs that Sofia presents? \newline A. Provide a set of images of the story, so she can recount it from beginning to end \newline B. Ask closed questions that aim at the recap of what was narrated, whose answers are yes or no \newline C. Offer paper and pencils so she can illustrate what was narrated, then questions oriented to the recount of the story will be asked \newline D. Include in a group of peers, in a dramatization activity of what was narrated, asking them to act out the story from beginning to end   \\
  \textbf{Gold}      & A       \\
  \textbf{Model}     & B  \\
  \textbf{Reasoning} & The user wants me to answer a multiple-choice question about curricular adaptations for a student named Sofia. \newline Let me analyze Sofia's profile: \newline - 2nd grade, 9 years old, special school \newline - Language and Communication subject \newline - Reading axis: identifies her written name and initial sound of words when they begin with vowels \newline - Oral communication axis: understands simple instructions and short texts read to her, verbally expresses herself with words like yes/no and words referring to basic needs (food, hygiene, health), accompanied by gestures \newline - Uses a communication board for oral expression \newline - Learning objective: Demonstrate comprehension of read narratives, extracting explicit and implicit information \newline - Evaluation: oral recount of the story \newline \textit{[\ldots\ truncated \ldots]} \newline - Evaluation (how we assess) \newline Here, the learning objective (comprehension) stays the same. The evaluation is adapted from open oral recount to closed yes/no questions. This respects the objective while accommodating her disability. \newline I will go with B.  \\
  \specialrule{1.2pt}{8pt}{8pt}
  \multicolumn{2}{@{}l}{\textbf{chile\_373}} \\
  \textbf{Question}  & A teacher is working with a middle school class on the unit: "The body as a reference for design" and wants to evaluate the objective: "Experiment with various languages, techniques, and modes of artistic and visual creation". Which of the following indicators should be present in the guideline that will evaluate the stated objective? \newline A. Draws sketches representing themes linked to the body as a support \newline B. Conducts bibliographic research on different concepts linked to the body as a support \newline C. Creates a prototype or clothing model that addresses the issue of the body as a support \newline D. Artistically investigates the materiality for the creation of a clothing prototype   \\
  \textbf{Gold}      & D       \\
  \textbf{Model}     & C  \\
  \textbf{Reasoning} & The user wants me to answer a multiple-choice question with ONLY the letter of the correct answer. No explanation. \newline The question is about evaluating an objective: "Experiment with various languages, techniques, and modes of artistic and visual creation" in the context of a unit "The body as a reference for design". \newline I need to determine which indicator should be present in the guideline to evaluate this specific objective. \newline Let's analyze each option: \newline A. "Draws sketches representing themes linked to the body as a support" - This involves drawing/sketching, which is a mode of artistic and visual creation. It fits "experiment with various languages, techniques, and modes". \newline B. "Conducts bibliographic research on different concepts linked to the body as a support" - This is research, not experimentation with artistic/visual creation techniques. It might relate to other objectives (like understanding concepts) but not specifically to experimenting with artistic languages and techniques. \newline C. "Creates a prototype or clothing model that addresses the issue of the body as a support" - This involves creating a prototype, which is a mode of artistic/visual creation and experimentation. This seems very aligned with the objective. \newline D. "Artistically investigates the materiality for the creation of a clothing prototype" - This focuses on investigating materiality, which is part of the process but more about investigation/research rather than the actual experimentation with languages/techniques/modes of creation itself. However, it does mention "artistically investigates". \newline Wait, I need to think about what makes the best indicator for "Experiment with various languages, techniques, and modes of artistic and visual creation". \newline \textit{[\ldots\ truncated \ldots]} \newline I think C is the strongest answer. \newline One more check: Could D be correct? "Artistically investigates the materiality for the creation of a clothing prototype". This focuses on materiality investigation. While materiality is important, the objective is about experimenting with languages, techniques, and modes of *creation*. Investigating materiality is preliminary to creation, not the creation itself. So D is less aligned than C. \newline Therefore, my final answer is C.  \\
  \specialrule{1.2pt}{8pt}{8pt}
  \multicolumn{2}{@{}l}{\textbf{chile\_993}} \\
  \textbf{Question}  & In the 5th grade of a school for deaf students, Macarena attends, who recently integrated into the course two years ago, as she previously attended a school with special educational needs. Although Macarena can now communicate fluently and adequately in Sign Language, she still faces challenges in writing, finding herself at a level well below the class average. Based on the above, which of the following strategies is most appropriate to promote the natural development of writing in Macarena? \newline A. Model the student's natural writing using a Communication Notebook \newline B. Reinforce writing through the creation of stories related to topics of interest to the student \newline C. Practice writing through the creation of texts supported by verb and vocabulary notebooks \newline D. Encourage writing based on sentences made in sign language by the student and written by the teacher   \\
  \textbf{Gold}      & A       \\
  \textbf{Model}     & D  \\
  \textbf{Reasoning} & The user wants me to answer a multiple-choice question about educational strategies for a deaf student named Macarena who is fluent in Sign Language but struggles with writing. \newline Let me analyze the scenario: \newline - Macarena is a 5th-grade deaf student \newline - She integrated into the school for deaf students two years ago \newline - Previously attended a school with special educational needs \newline - Communicates fluently in Sign Language \newline - Faces challenges in writing, well below class average \newline The question asks which strategy is most appropriate to promote the NATURAL development of writing in Macarena. \newline Let's evaluate each option: \newline \textit{[\ldots\ truncated \ldots]} \newline The principle is: for deaf students, sign language is their natural language, and written Spanish is a second language. Therefore, literacy should be taught through sign language. \newline Option D describes exactly this: the student produces in sign language, and this is written down. This helps her see her language in written form. \newline I'll go with D.  \\
  \specialrule{1.2pt}{8pt}{8pt}
  \multicolumn{2}{@{}l}{\textbf{chile\_624}} \\
  \textbf{Question}  & Read the following situation: \newline At the start of a thematic unit, a teacher explains the concept of habituation through a definition. In this context, a student asks for clarification through an example. \newline Which of the following statements represents an appropriate example to explain the concept of habituation? \newline A. If we repeatedly put ice on our hand, it will become numb and we will not feel the ice as cold \newline B. We generally buy the same shampoo because we have become accustomed to it \newline C. By lying repeatedly, we generate distrust in others \newline D. If a car parked on the street sounds its alarm every night at 2 a.m., at some point we will stop caring   \\
  \textbf{Gold}      & A       \\
  \textbf{Model}     & D  \\
  \textbf{Reasoning} & The user wants me to answer a multiple-choice question about the concept of habituation. I need to choose the correct answer letter only, with no explanation. \newline Let me analyze the concept of habituation first. Habituation is a psychological learning process where there is a decrease in response to a stimulus after being repeatedly exposed to it. It's a form of non-associative learning. The key is that the stimulus itself is neither harmful nor rewarding, and the organism simply stops responding to it because it has become familiar or irrelevant. \newline Now let's evaluate each option: \newline A. "If we repeatedly put ice on our hand, it will become numb and we will not feel the ice as cold" - This describes sensory adaptation or neural adaptation (numbness due to physical changes in receptors/neurons), not habituation. Habituation is a behavioral/psychological decrease in response, not a physiological numbing. \newline B. "We generally buy the same shampoo because we have become accustomed to it" - This is more about preference, routine, or possibly positive reinforcement (if the shampoo works well). It's not habituation. Habituation isn't about consumer loyalty. \newline C. "By lying repeatedly, we generate distrust in others" - This is about social consequences of behavior, not habituation. \newline D. "If a car parked on the street sounds its alarm every night at 2 a.m., at some point we will stop caring" - This is the classic example of habituation. The alarm is a stimulus that initially elicits a response (annoyance, awakening, concern), but after repeated exposure without negative consequence, the response diminishes. We stop reacting to it. This perfectly illustrates habituation. \newline So the correct answer is D.  \\
  \specialrule{1.2pt}{8pt}{8pt}
  \multicolumn{2}{@{}l}{\textbf{chile\_660}} \\
  \textbf{Question}  & A high school teacher is preparing a class on Kantian theory about the possibility of knowledge; to do this, they must determine the prior learning experiences necessary for their students to adequately address such a theory. What prior learning should the teacher recognize in their students to tackle this theory? \newline A. Recognize the ideas of D. Hume and J. Locke as theoretical antecedents to the proposition of I. Kant \newline B. Understand English dogmatism and empiricism as currents preceding the Kantian proposal \newline C. Identify the notions of time and space as structures of thought necessary to know, according to I. Kant \newline D. Understand the concepts of phenomenon and noumenon as Kantian explanations about what can be known   \\
  \textbf{Gold}      & B       \\
  \textbf{Model}     & A  \\
  \textbf{Reasoning} & The user wants me to answer a multiple-choice question about Kantian theory and prior learning experiences necessary for students to adequately address such a theory. \newline Let me analyze the question carefully. The teacher is preparing a class on Kantian theory about the possibility of knowledge. They need to determine prior learning experiences necessary for students to adequately address such a theory. The question asks what prior learning the teacher should recognize in their students. \newline Let's look at each option: \newline A. Recognize the ideas of D. Hume and J. Locke as theoretical antecedents to the proposition of I. Kant \newline - This refers to prior learning about Hume and Locke, who were indeed major influences on Kant. Hume's skepticism awakened Kant from his "dogmatic slumber," and Locke's empiricism was part of the context. This seems like legitimate prior knowledge that would help students understand Kant's critical philosophy, which was partly a response to rationalism and empiricism. \newline B. Understand English dogmatism and empiricism as currents preceding the Kantian proposal \newline - This is somewhat similar to A but less precise. "English dogmatism" is not a standard term; usually, it's rationalism (Continental, e.g., Descartes, Leibniz) vs. empiricism (British/English, e.g., Locke, Hume, Berkeley). This option seems historically inaccurate or at least imprecise because dogmatism/rationalism is more associated with Continental philosophers, not specifically English. Also, empiricism is English/British, but dogmatism isn't typically characterized as English. \newline C. Identify the notions of time and space as structures of thought necessary to know, according to I. Kant \newline - This is actually part of Kant's theory itself (the Transcendental Aesthetic), not prior learning. If students already know this, they already know part of what is to be taught. The question asks for prior learning necessary to tackle the theory, not what constitutes the theory itself. While having some familiarity might help, this seems like it would be part of the new content, not prior learning. \newline \textit{[\ldots\ truncated \ldots]} \newline Thus, the answer is A. \newline Actually, let me think once more. Could "dogmatism" refer to something else? In Kant, "dogmatism" is the opposite of skepticism and criticism. He criticizes dogmatic metaphysics. But the dogmatic metaphysics he primarily had in mind was the Leibnizian-Wolffian school in Germany. So "English dogmatism" is definitely wrong. \newline Therefore, A is the correct answer.  \\
  \specialrule{1.2pt}{8pt}{8pt}
  \multicolumn{2}{@{}l}{\textbf{chile\_666}} \\
  \textbf{Question}  & A teacher will begin working on the unit 'Economic models and their impact on the territory', with the goal of having her high school students analyze the transformations experienced by the geographical space as a result of the relationships between economic, demographic, and spatial processes. Therefore, she has decided to start the class by recalling the meaning of the term geographical space, through a group brainstorming session. In the responses given by the students, those that characterized the geographical space exclusively from its physical or natural component prevailed, attributing to it an attribute of stability in the medium and long term. Considering the didactic principles that guide the teaching of geography, which one of them is essential to work on in order to expand the concept of space that this group of students has? \newline A. Connection \newline B. Location \newline C. Interdependence \newline D. Causality   \\
  \textbf{Gold}      & C       \\
  \textbf{Model}     & A  \\
  \textbf{Reasoning} & The user wants me to answer a multiple-choice question about didactic principles in geography teaching. \newline Let me analyze the question: \newline - Unit: 'Economic models and their impact on the territory' \newline - Goal: analyze transformations experienced by geographical space as a result of relationships between economic, demographic, and spatial processes \newline - Current student understanding: geographical space characterized exclusively from its physical/natural component, attributing stability in medium/long term \newline - Need: expand the concept of space beyond just physical/natural stability \newline The students see geographical space as just physical/natural and stable. But the teacher wants them to understand it as something that transforms due to economic, demographic, and spatial processes. \newline Which didactic principle is essential to work on? \newline Let's look at the options: \newline \textit{[\ldots\ truncated \ldots]} \newline - Causalidad \newline And the principle of connection is indeed the one that integrates the components of geographical space. \newline Final answer: A.  \\
  \specialrule{1.2pt}{8pt}{8pt}
  \multicolumn{2}{@{}l}{\textbf{chile\_610}} \\
  \textbf{Question}  & A teacher in a Third-Year class was working with the students on the exponential function and in the next class will work on the logarithmic function. To do this, she will carry out a guided activity in which, starting from the algebraic expression f(x) = log10(x), they will tabulate some values for x with their evaluation in this function and will sketch its graph, determining its domain and range through an algebraic analysis. For this reason, the teacher decides to complement the algebraic analysis using another methodological strategy. Considering this pedagogical purpose, which of the following strategies carried out by the teacher promotes the understanding of the domain and range of logarithmic functions? \newline A. Pose situations modeled by logarithmic functions expressed algebraically, interpreting, together, their domain and range concerning the given context \newline B. Project the graphs of different logarithmic functions, each one with its respective inverse exponential function, determining the domain and range of the first, based on the second \newline C. Pose the algebraic expression of different logarithmic functions, asking the students to obtain their corresponding domain and range, expressed in mathematical notation \newline D. Project the graphs of different logarithmic functions, each one with its respective algebraic expression, analyzing, together, the differences and similarities observed in their domain and range   \\
  \textbf{Gold}      & D       \\
  \textbf{Model}     & B  \\
  \textbf{Reasoning} & The user wants me to answer a multiple-choice question about teaching strategies for logarithmic functions, specifically regarding domain and range. \newline Let me analyze the scenario: \newline - Teacher worked on exponential functions \newline - Next class will work on logarithmic functions \newline - Students will start with f(x) = log10(x), tabulate values, sketch graph, determine domain and range through algebraic analysis \newline - Teacher wants to complement the algebraic analysis using another methodological strategy \newline - Goal: promote understanding of domain and range of logarithmic functions \newline Now let's evaluate each option: \newline A. Pose situations modeled by logarithmic functions expressed algebraically, interpreting, together, their domain and range concerning the given context \newline \textit{[\ldots\ truncated \ldots]} \newline Another strategy: could be graphical, contextual, or using inverse relationships. \newline In mathematics education, when teaching logarithms after exponentials, the most common and recommended strategy is to exploit the inverse relationship. This is exactly what option B describes. \newline Therefore, the answer is B.  \\
  \specialrule{1.2pt}{8pt}{8pt}
  \multicolumn{2}{@{}l}{\textbf{chile\_574}} \\
  \textbf{Question}  & A fourth-grade teacher is planning a class in which she will work on comparing decimal numbers. To start, she will ask the students to compare the following numbers: 1.37 with 1.4 and 1.25 with 1.5; for this purpose, she is looking for a type of representation that will help them perform the exercise and understand the content. Which of the following activities would meet the requirement that the teacher has set? \newline A. Write the quantities as hundredths \newline B. Express the quantities in fractional form \newline C. Represent the quantities with concrete materials \newline D. Illustrate the quantities in everyday situations   \\
  \textbf{Gold}      & A       \\
  \textbf{Model}     & C  \\
  \textbf{Reasoning} & The user wants me to answer a multiple-choice question about teaching decimal comparison to fourth-grade students. I need to choose the best activity that helps students compare decimals like 1.37 vs 1.4 and 1.25 vs 1.5. \newline Let me analyze each option: \newline A. Write the quantities as hundredths - This would mean writing 1.4 as 1.40 and 1.5 as 1.50. This is a good mathematical strategy for comparing decimals, but the question asks for a "type of representation" that helps them perform the exercise and understand the content. Writing as hundredths is more of a procedural step. \newline B. Express the quantities in fractional form - Converting to fractions (like 137/100 vs 14/10 or 140/100) might actually make it harder for fourth graders to compare, and it's not the most intuitive representation for understanding decimal magnitude. \newline C. Represent the quantities with concrete materials - Using base-ten blocks, decimal squares, number lines, or other manipulatives would provide a concrete representation. This is often recommended in math education for conceptual understanding. For example, showing that 1.4 is 1 whole and 4 tenths, while 1.37 is 1 whole and 3 tenths and 7 hundredths - visually seeing that 4 tenths \textgreater{} 3 tenths and 7 hundredths. \newline D. Illustrate the quantities in everyday situations - Real-world contexts (like money) could help, but money might be confusing since \$1.4 would be \$1.40, and students might not naturally think of \$1.4 as valid notation. Everyday situations are good for application but may not be the best "representation" for the initial comparison task. \newline Wait, let me reconsider. The question specifically mentions "a type of representation." In mathematics education, especially for decimals, concrete materials (manipulatives) are considered representations. But writing as hundredths is also a representational strategy. \newline Actually, looking at typical pedagogical approaches for teaching decimal comparison: \newline - A common misconception is that 1.37 \textgreater{} 1.4 because 37 \textgreater{} 4. \newline \textit{[\ldots\ truncated \ldots]} \newline If they just write 1.40, they might understand that 140 \textgreater{} 137, but they might still not understand WHY 1.4 is the same as 1.40 conceptually. \newline Therefore, C seems pedagogically superior for understanding. \newline I'm going to go with C.  \\
  \specialrule{1.2pt}{8pt}{8pt}
  \multicolumn{2}{@{}l}{\textbf{chile\_714}} \\
  \textbf{Question}  & A teacher will start working with her 6th grade students based on the following plan: 6th Grade Unit: The political organization of Chile and the basic principles established by the Constitution. Learning Objective: Explain some fundamental elements of the democratic organization of Chile, including the division of state powers and their representation through elected officials (councilors, mayors, deputies, senators, and President). Evaluation indicator: Recognize the main characteristics that a representative democratic regime must have. Which of the following contents is essential for the students to manage in order to achieve the expected performance? \newline A. Citizen participation \newline B. Rights and duties \newline C. Type of State \newline D. Method of appointing authorities   \\
  \textbf{Gold}      & A       \\
  \textbf{Model}     & D  \\
  \textbf{Reasoning} & The user wants me to answer a multiple-choice question about a 6th grade unit on the political organization of Chile and democratic principles. \newline Let me analyze the question: \newline Context: \newline - Unit: The political organization of Chile and the basic principles established by the Constitution. \newline - Learning Objective: Explain some fundamental elements of the democratic organization of Chile, including the division of state powers and their representation through elected officials (councilors, mayors, deputies, senators, and President). \newline - Evaluation indicator: Recognize the main characteristics that a representative democratic regime must have. \newline Question: Which of the following contents is essential for the students to manage in order to achieve the expected performance? \newline Options: \newline A. Citizen participation \newline \textit{[\ldots\ truncated \ldots]} \newline But let me think if there's any trick. "Citizen participation" could be argued as essential to democracy. However, the specific focus on "elected officials" and "representative democratic regime" points more specifically to the method of appointing authorities (election/representation) rather than participation in general. \newline Also, in many standardized tests for teacher evaluation in Chile (like the PAES or teacher certification exams), this type of question tests alignment between objectives and content. The content that most directly supports the objective about elected officials and representative democracy would be the method of appointing authorities. \newline I'll go with D.  \\
  \specialrule{1.2pt}{8pt}{8pt}
  \multicolumn{2}{@{}l}{\textbf{chile\_960}} \\
  \textbf{Question}  & Read the following situation: The classroom team of a 2nd grade, in the subject of Language and Communication, is carrying out the following activity: 1. Shared reading 2. Questions to the students 3. Emilio's response The Monkey Tito Tito is not an ordinary monkey. Tito does not like to climb trees and hates to eat bananas. He prefers to stroll, smell the flowers, and collect the nuts that fall from the trees. According to what has been presented, what is the support requirement that Emilio presents? \newline A. Locate information about the characters in a text \newline B. Locate details about the events in a text \newline C. Draw conclusions from explicit information in a text \newline D. Establish time and place relationships between the actions in a text   \\
  \textbf{Gold}      & C       \\
  \textbf{Model}     & A  \\
  \textbf{Reasoning} & The user wants me to answer a multiple-choice question based on a reading comprehension scenario. \newline Let me analyze the situation: \newline - 2nd grade Language and Communication class \newline - Activity: Shared reading Questions to students Emilio's response \newline - Text: "The Monkey Tito" - describes Tito as not ordinary, doesn't like climbing trees or eating bananas, prefers strolling, smelling flowers, collecting nuts. \newline The question asks: What is the support requirement that Emilio presents? \newline Wait, the prompt doesn't actually tell us what Emilio's response was. It just says "Emilio's response" as step 3, but doesn't provide the content of his response. This seems like missing information. \newline However, looking at the options: \newline A. Locate information about the characters in a text \newline \textit{[\ldots\ truncated \ldots]} \newline - Drawing conclusions requires higher-order thinking. \newline The question asks for the "support requirement" (necesidad de apoyo). Usually, support is needed for the specific skill being assessed. Since the text is a character description, the skill is locating character information. \newline I'll go with A.  \\
  \specialrule{1.2pt}{8pt}{8pt}
  \multicolumn{2}{@{}l}{\textbf{chile\_787}} \\
  \textbf{Question}  & The technology and science teachers of a 5th grade have agreed to work in parallel on the following science topic: the growth of a plant. For this, the technology teacher has asked her students to develop an explanatory presentation of this process, in order to reinforce the knowledge acquired in science and, additionally, to create an educational material for the 3rd grade students, who also need to learn this process. \newline Which of the following digital tools is the most appropriate for developing the presentation, considering the purpose and the level of the students? \newline A. PowerPoint, using text and images \newline B. Movie Maker, recording and editing a video \newline C. AnimatorDV, to make an animation of the process \newline D. Padlet, publishing useful content to understand the process   \\
  \textbf{Gold}      & C       \\
  \textbf{Model}     & A  \\
  \textbf{Reasoning} & The user wants me to answer a multiple-choice question with ONLY the letter of the correct answer. No explanation. \newline Let me analyze the question: \newline - Context: 5th grade students creating an explanatory presentation about plant growth. \newline - Purpose: Reinforce knowledge acquired in science AND create educational material for 3rd grade students. \newline - Need to choose the most appropriate digital tool considering purpose and student level. \newline Options: \newline A. PowerPoint, using text and images \newline B. Movie Maker, recording and editing a video \newline C. AnimatorDV, to make an animation of the process \newline \textit{[\ldots\ truncated \ldots]} \newline So it's between A and C. Given that the teacher asked for an "explanatory presentation," and PowerPoint is the presentation tool par excellence, and considering that 5th graders can handle it well while 3rd graders can understand it, I think A is the intended answer. \newline Actually, one more thought: In some educational frameworks, having 5th graders create materials for 3rd graders is an example of "peer tutoring" or "learning by teaching." The tool should allow clear explanation of a process. PowerPoint's linear slide structure helps organize the stages of plant growth sequentially. \newline I'll go with A.  \\
  \specialrule{1.2pt}{8pt}{8pt}
  \multicolumn{2}{@{}l}{\textbf{chile\_128}} \\
  \textbf{Question}  & A teacher of a preschool class is developing the unit Thinking like a scientist. The children have actively participated in pedagogical experiences such as conducting experiments and learning about the scientific method. The skills they will practice are observation and inference. Which pedagogical resource would be suitable to promote the development of these skills in the children? \newline A. Metacognition sheet, to record the new concepts learned during the experiment and then draw conclusions \newline B. Tracking guide, to record the initial ideas and observations of the experiment, and then compare the data obtained afterwards \newline C. Planning log, through which the steps and initial ideas of the experiment are organized, and then the experience is monitored \newline D. Collection and observation tools that facilitate and allow a more precise exploration, to then describe the findings obtained   \\
  \textbf{Gold}      & B       \\
  \textbf{Model}     & D  \\
  \textbf{Reasoning} & The user wants me to answer a multiple-choice question with ONLY the letter of the correct answer. No explanation. \newline Let me analyze the question: \newline - Context: Preschool class, unit "Thinking like a scientist" \newline - Children have participated in experiments and learned about the scientific method \newline - Skills to practice: observation and inference \newline - Question: Which pedagogical resource is suitable to promote these skills? \newline Let's evaluate each option: \newline A. Metacognition sheet, to record new concepts learned and draw conclusions \newline - Metacognition is thinking about thinking. While drawing conclusions is good, this focuses more on recording concepts and conclusions rather than the process of observation and inference itself. Also, "metacognition sheet" might be too abstract for preschoolers. \newline \textit{[\ldots\ truncated \ldots]} \newline Wait, I should double-check. Is there any reason B would be better? "Tracking guide" - in some curricula, they use observation tracking guides. But for preschoolers who may not read/write well, tools are better than guides/logs/sheets. \newline The question asks for a "pedagogical resource." All options are resources. But D provides the concrete means to observe, which is fundamental. You can't infer without good observation, and tools make observation more precise. \newline Yes, I'll go with D.  \\
  \specialrule{1.2pt}{8pt}{8pt}
  \multicolumn{2}{@{}l}{\textbf{chile\_1026}} \\
  \textbf{Question}  & Read the following situation: Asuncian is a four-year-old student who has recently been diagnosed with Autism Spectrum Disorder and has entered a special school. She shows challenges in interpersonal relationships; she responds appropriately to some actions proposed to her, but she is inflexible in the face of changes in spaces and routines. Her mother reports that it is very difficult to leave home with Asuncian to visit someone, go to the park, or shop, as she constantly cries, so her outings have been reduced only to taking her to school. The teacher and the psychologist at the institution suggest to Asuncian's mother some actions to follow: Visual (photography) and verbal (short) anticipation before and after outings. Positive reinforcement when the student shows the slightest expected behavior. Considering Asuncian's characteristics, which of the following options most adequately justifies the relevance of the anticipation actions proposed by the professionals at this establishment? \newline A. They will provide her with security through social reinforcement \newline B. They will give her security through academic reinforcement \newline C. They will provide her with security through affective reinforcement \newline D. They will provide her with security through practicing her cognitive skills   \\
  \textbf{Gold}      & A       \\
  \textbf{Model}     & D  \\
  \textbf{Reasoning} & The user wants me to answer a multiple-choice question based on a case study about a child with Autism Spectrum Disorder. \newline Let me analyze the situation: \newline - Asuncian is 4 years old, diagnosed with ASD \newline - She has challenges in interpersonal relationships \newline - She is inflexible with changes in spaces and routines \newline - Going out is difficult, she cries, outings reduced to only school \newline - Professionals suggest: visual (photography) and verbal (short) anticipation before and after outings, plus positive reinforcement for slightest expected behavior. \newline The question asks: Which option most adequately justifies the relevance of the anticipation actions proposed? \newline Anticipation strategies (visual and verbal) for children with ASD are typically used to help them understand what will happen next, reducing anxiety about unpredictability and changes in routine. This provides a sense of security because they can predict the sequence of events. \newline \textit{[\ldots\ truncated \ldots]} \newline These are cognitive skills. The security comes from this cognitive practice. \newline I think D is correct. \newline Final answer: D  \\
  \specialrule{1.2pt}{8pt}{8pt}
  \multicolumn{2}{@{}l}{\textbf{chile\_1017}} \\
  \textbf{Question}  & Read the following situation: Martina is a four-year-old student attending first level of transition. She began to verbalize her first words around two and a half years, mainly to satisfy immediate needs. Currently, she maintains a peculiar jargon, which is sometimes characterized as unintelligible, even though her pronunciation is perfect when she verbalizes words she has heard on television. She also repeats some phrases from cartoons that she likes, imitating prosodic patterns that include the production of stress and tone of the heard pattern. The characteristics of Martina's language have drawn the attention of her teacher, who wishes to conduct a psychopedagogical assessment to determine the origin of these characteristics, as she suspects the presence of an autism spectrum disorder in the student. Which of the following procedures should the teacher carry out to conduct the psychopedagogical assessment of Martina? \newline A. Conduct an inventory of behaviors and skills aimed at identifying Martina's learning patterns and abilities in order to subsequently create guidelines and work strategies \newline B. Provide a guideline with a rating scale to the family in which the different areas of development are evaluated, and unify the resulting criteria from this evaluation with observations in the classroom \newline C. Request a psychometric assessment from the family to determine Martina's IQ, and based on the results, create a work structure that references those areas of development that present challenges \newline D. Prepare a checklist that evaluates the different objectives proposed in the Early Childhood Education Curriculum Framework to determine the areas that require support and define those that are satisfactorily achieved   \\
  \textbf{Gold}      & A       \\
  \textbf{Model}     & B  \\
  \textbf{Reasoning} & The user wants me to answer a multiple-choice question about a psychopedagogical assessment for a 4-year-old student named Martina who shows signs of autism spectrum disorder (ASD). \newline Let me analyze the situation: \newline - Martina is 4 years old, in first level of transition (preschool/kindergarten) \newline - Delayed first words (2.5 years) \newline - Uses peculiar jargon, sometimes unintelligible \newline - Perfect pronunciation when repeating words from TV \newline - Repeats phrases from cartoons (echolalia) \newline - Imitates prosodic patterns including stress and tone \newline - Teacher suspects autism spectrum disorder \newline \textit{[\ldots\ truncated \ldots]} \newline Option C focuses on IQ which is wrong. \newline Option D focuses only on curriculum objectives, missing the developmental/behavioral aspect. \newline Therefore, the answer is B.  \\
  \specialrule{1.2pt}{8pt}{8pt}
  \multicolumn{2}{@{}l}{\textbf{chile\_460}} \\
  \textbf{Question}  & To start the unit Design and Dissemination, in a 2nd Year secondary class, a teacher conducts a diagnostic assessment to determine the students' approach to design and its different branches. Once the assessment is applied, the teacher identifies that they have a confused notion about industrial design, as they exclusively relate it to the manufacturing of industrial parts on a large scale. What activity could the teacher propose to the students so they can understand the concept of industrial design? \newline A. Observe pieces of industrial design to analyze them, classify them, and make sketches of the exhibited pieces \newline B. Choose design pieces found in the classroom to investigate their origin, evolution, and relationship with their production context \newline C. Deconstruct an industrial design piece through photography, for the purpose of creating a new design piece that addresses current needs \newline D. Research iconic design pieces present in the history of industrial design, to subsequently create a design piece considering contextual needs   \\
  \textbf{Gold}      & B       \\
  \textbf{Model}     & D  \\
  \textbf{Reasoning} & The user wants me to answer a multiple-choice question with ONLY the letter (A, B, C, or D). No explanation. \newline The scenario: 2nd Year secondary class, starting a unit on Design and Dissemination. Diagnostic assessment shows students think industrial design is exclusively about manufacturing industrial parts on a large scale. Teacher wants to propose an activity so they can understand the concept of industrial design. \newline I need to choose the best activity to correct this misconception. The misconception is that industrial design = manufacturing industrial parts at large scale. Industrial design is actually about designing products for mass production, considering aesthetics, functionality, ergonomics, etc. It's not just about mechanical/industrial parts, but everyday objects, furniture, appliances, etc. \newline Let's analyze each option: \newline A. Observe pieces of industrial design to analyze them, classify them, and make sketches of the exhibited pieces. \newline - This involves observation and classification. It might help them see what industrial design encompasses. But does it directly address the misconception about manufacturing parts? It shows them examples, which could broaden their view. However, it's somewhat passive. \newline B. Choose design pieces found in the classroom to investigate their origin, evolution, and relationship with their production context. \newline - This connects everyday objects (found in classroom) to their design and production context. It helps students realize that industrial design includes everyday objects, not just heavy industrial parts. Investigating origin, evolution, and production context directly addresses the context of mass production and design. This seems strong because it uses familiar objects. \newline C. Deconstruct an industrial design piece through photography, for the purpose of creating a new design piece that addresses current needs. \newline \textit{[\ldots\ truncated \ldots]} \newline D: research, create. (Researching history gives conceptual foundation; creating applies it.) \newline I think D is the intended answer because it explicitly mentions "history of industrial design" which is the standard way to teach the breadth of the concept, and "considering contextual needs" captures the human-centered aspect that distinguishes it from mere manufacturing. \newline Final decision: D.  \\
  \specialrule{1.2pt}{8pt}{8pt}
  \multicolumn{2}{@{}l}{\textbf{chile\_726}} \\
  \textbf{Question}  & A teacher is planning a class for 5th grade with the following learning objective: "Characterize the landscapes of the large natural regions of Chile." Which of the following activities allows addressing the proposed objective? \newline A. Construct a model of the characteristic landscape of a natural zone highlighting the most well-known places, to exhibit it in the school yard \newline B. Research on the site www.sernatur.cl about the natural resources and economic activities associated with each zone, to present them in a diagram \newline C. Interview acquaintances who have been in another geographic zone, asking them about the landscapes and customs they saw, to create a travel journal \newline D. Create a tourist brochure with images of the landscapes promoting the geographical advantages and disadvantages of visiting each zone, to set up a poster exhibition   \\
  \textbf{Gold}      & D       \\
  \textbf{Model}     & A  \\
  \textbf{Reasoning} & The user wants me to evaluate a multiple-choice question and answer with ONLY the letter of the correct answer. \newline Learning objective: "Characterize the landscapes of the large natural regions of Chile." \newline Let's analyze each option: \newline A. Construct a model of the characteristic landscape of a natural zone highlighting the most well-known places, to exhibit it in the school yard. \newline - This involves characterizing landscapes (physical representation, highlighting features). It directly addresses characterizing landscapes of natural regions. \newline B. Research on the site www.sernatur.cl about the natural resources and economic activities associated with each zone, to present them in a diagram. \newline - This focuses on natural resources and economic activities, not specifically on characterizing landscapes. \newline C. Interview acquaintances who have been in another geographic zone, asking them about the landscapes and customs they saw, to create a travel journal. \newline - This includes landscapes but also customs (cultural aspect). The objective is specifically about landscapes of large natural regions. While it touches on landscapes, it relies on second-hand accounts and mixes in customs which isn't the focus. Also, "large natural regions of Chile" - interviewing acquaintances might not cover all regions systematically. \newline \textit{[\ldots\ truncated \ldots]} \newline I think A is the intended answer. \newline Final check: The user wants ONLY the letter. No explanation. \newline Answer: A  \\
  \specialrule{1.2pt}{8pt}{8pt}
  \multicolumn{2}{@{}l}{\textbf{chile\_683}} \\
  \textbf{Question}  & A teacher tells his 3rd year high school students that during the unit "The Democratic Breakdown", the following objective will be worked on in groups: "Analyze the breakdown of democracy from a multi-causal perspective". Therefore, he proposes that they search for information about the period in different sources. Then, he explains that based on this information, in groups, they will share what they have investigated with the rest of the class. One of the groups asks the teacher if it is okay to make their presentation with a timeline of the 1970-1973 period, highlighting relevant and diverse historical events. The teacher, after hearing them, invites them to look for another form of presentation, as the proposed activity would not allow them to fulfill the selected learning objective. Why is the timeline resource inappropriate with respect to the learning objective proposed to the students? \newline A. Because it favors analysis from a chronological perspective, rather than from multiple perspectives \newline B. Because it distinguishes historical events, rather than establishing a hierarchy according to their importance \newline C. Because it limits the analysis to a linear view of events, rather than an understanding of the historical process \newline D. Because it makes it easy to order events, rather than grouping them according to the different dimensions involved in the historical process   \\
  \textbf{Gold}      & C       \\
  \textbf{Model}     & D  \\
  \textbf{Reasoning} & The user wants me to answer a multiple-choice question with ONLY the letter of the correct answer. No explanation. \newline Let me analyze the question: \newline Context: A teacher asks students to "Analyze the breakdown of democracy from a multi-causal perspective" in groups. One group wants to make a timeline of 1970-1973 highlighting relevant and diverse historical events. The teacher says this won't allow them to fulfill the learning objective. Why is the timeline inappropriate? \newline The key is the learning objective: "multi-causal perspective" (perspectiva multicausal). This means analyzing multiple causes, dimensions, factors (political, economic, social, international, etc.) that led to the breakdown of democracy. \newline A timeline is chronological. It orders events sequentially. While it can show diverse events, its primary structure is temporal/chronological. The teacher wants multi-causal analysis. \newline Let's look at the options: \newline A. Because it favors analysis from a chronological perspective, rather than from multiple perspectives \newline - This says timeline favors chronological over multiple perspectives. But "multiple perspectives" could mean different viewpoints (e.g., different social groups), not necessarily multi-causal. However, in Spanish educational context, "perspectiva multicausal" is about causes. But "multiple perspectives" is close. \newline B. Because it distinguishes historical events, rather than establishing a hierarchy according to their importance \newline \textit{[\ldots\ truncated \ldots]} \newline But wait, I should consider if there's any trick. The group wants to highlight "relevant and diverse historical events." The teacher says the timeline won't allow them to fulfill the objective. Why? Because a timeline organizes by time, not by cause. To analyze multi-causal, you need to see relationships, influences, dimensions. A timeline is sequential. \newline Option D is the most accurate. \newline I will answer D.  \\
  \specialrule{1.2pt}{8pt}{8pt}
  \multicolumn{2}{@{}l}{\textbf{chile\_1040}} \\
  \textbf{Question}  & Read the following situation: Agustin's teacher, a student with Autism Spectrum disorder in 1st grade, is working with the class on deepening positional value and mathematical reasoning, in relation to everyday contexts in which students use numbers. To do this, he wants to propose problems to the class involving counting in increments of 10 and 50, up to one hundred. For example, they must buy a fruit for a maximum price of 100 pesos and give or ask for change in coins of 10 and 50 pesos. To consolidate the understanding of this counting, prior to resolving problems that involve addition, the teacher invites the students to carry out the following exercise: Place in the box the coins needed to complete 100 pesos: Based on the above, why does the preparation strategy favor Agustin's learning? \newline A. It allows him to understand the new content by approaching it through visual aids \newline B. It allows him to gradually anticipate more complex mathematical exercises and avoid setbacks when facing the unknown \newline C. It allows him to internalize the reviewed content through the repetition of an exercise he will have to apply in problem solving \newline D. It allows him to put into practice his knowledge about transactions with coins in understanding mathematical concepts   \\
  \textbf{Gold}      & A       \\
  \textbf{Model}     & B  \\
  \textbf{Reasoning} & The user wants me to answer a multiple-choice question based on a situation involving a student with Autism Spectrum Disorder (Agustin) in 1st grade. \newline Let me analyze the situation: \newline - Agustin has Autism Spectrum Disorder \newline - The class is working on deepening positional value and mathematical reasoning \newline - Context: everyday use of numbers \newline - Problems involve counting in increments of 10 and 50, up to 100 \newline - Example: buying fruit for max 100 pesos, giving/asking for change in coins of 10 and 50 \newline - Prior to addition problems, the teacher invites students to do an exercise: "Place in the box the coins needed to complete 100 pesos" \newline - Question: Why does this preparation strategy favor Agustin's learning? \newline \textit{[\ldots\ truncated \ldots]} \newline The question says "why does the preparation strategy favor Agustin's learning?" The term "preparation strategy" in special education for ASD is often about preparing the student for upcoming tasks, using visual schedules, advance organizers, etc. The purpose is to make the unknown known. \newline However, looking at the options, B says "avoid setbacks when facing the unknown." This is a very strong match for ASD characteristics (rigidity, fear of unknown, need for predictability). The preparation gives him a known framework before the complex problem. \newline I think B is the intended answer.  \\
  \specialrule{1.2pt}{8pt}{8pt}
  \bottomrule
\end{longtable}